\section{Related Work}%
\label{sec:related-work}

Most state-of-the-art approaches~\citep{Lehmann2009DL-Learner,Lehmann2010Concept,Tran2012Approach} for class expression learning employ inductive logic programming (ILP) with \emph{refinement operators}. The first ILP approaches GOLEM~\citep{Muggleton1990Efficient,Muggleton1991Inductive}, ProGOLEM~\citep{Muggleton2009ProGolem}, and Aleph~\citep{Muggleton1995Inverse} employed \emph{upward} refinement operators by means of least generalization. However, as \citet{Badea2000Refinement} argue, they generate overly specific concepts that tend to overfit. Hence, they propose tackling the problem by \emph{downward} refinement operators with subsequent approaches improving the downward refinement operators by means of heuristics~\citep{Lehmann2010Concept,Lehmann2011Class} and parallelization~\citep{Tran2012Approach}. The latest approaches, DL-FOIL~\citep{Fanizzi2018DLFoil} and SPaCEL~\citep{Tran2017Parallel}, employ a hybrid approach of \emph{upward} and \emph{downward} refinements. DL-FOIL, however, assumes the existence of a ``perfect'' concept covering exactly all positive and negative instances, which is hardly possible in realistic scenarios and causes the algorithm not to terminate. SPaCEL overcomes this problem by combining many partial descriptions. Our approach, Evo\-Learner, finds shorter concepts that are more likely to generalize, and in cases without a perfect solution, our approach terminates with a good approximation. In our evaluation, we show that  EvoLearner outperforms SPaCEL on most datasets.

While most ILP approaches employ refinement operators, there have been some attempts in the direction of \emph{evolutionary algorithms}. \mbox{\citet{Reiser1999Evolution}} showed that the ability of evolutionary approaches to search in a global space, allowing to escape local minima, can lead to an improvement over ILP algorithms. To our knowledge, \citet{Lehmann2007Hybrid} was the first to apply evolutionary algorithms to the task of concept learning in \emph{description logic}. He combined standard genetic programming~(GP) approaches with genetic refinement operators. For evaluation, he employs a single, small dataset comparing his genetic refinement approach to a standard GP approach. However, we were not able to reproduce his reported results, and in a pilot study, we found that our evolutionary approach works better with standard GP operators than with genetic refinements. \citet{Divina2006Evolutionary} gives an overview of evolutionary concept learning in first-order logic. In contrast, we specialize in description logics and contribute a sophisticated initialization strategy based on random walks~\citep{Fronczak2009Biased, Sinatra2011Maximal}.
